% TOP_PROBLEM: {"id": 11300005, "title": "Nearby Lagrangian conjecture", "category_id": 7, "category": {"id": 7, "name": "topology", "display_name": "Topology", "description": "Properties preserved under continuous deformations.", "slug": "topology", "order_index": 7, "created_at": "2026-07-31T15:26:25.671Z"}, "status": "partially_solved"}
% FIRST_POSTED: 2026-10-06
% ATTEMPT_STATUS: solved
% OPENAI_CLAIM: {"schema_version":1,"source_repo":"https://github.com/openai/math","source_commit":"adc7f1241b42e322a6451854ab7e4b4c146bf78a","release_date":"2026-10-06","match":"full","resolution":"disproved","independently_reviewed":false,"second_pass":"agreed","families":[{"family":"340","title":"A counterexample to the nearby Lagrangian conjecture","subject":"Differential geometry","manuscripts":[{"title":"A counterexample to the nearby Lagrangian conjecture","pdf_path":"preprints/A-counterexample-to-the-nearby-Lagrangian-conjecture-September-23-2026/paper.pdf","readme_path":"preprints/A-counterexample-to-the-nearby-Lagrangian-conjecture-September-23-2026/README.md","date":"2026-09-23","theorem_ref":"thm:main (sections/introduction.tex:25)","statement_sources":[{"path":"preprints/A-counterexample-to-the-nearby-Lagrangian-conjecture-September-23-2026/build/sections/introduction.tex","line":25,"label":"thm:main"}]}],"lean":null,"reasoning_trace":null}]}
\documentclass[11pt]{article}
\usepackage[margin=25mm]{geometry}
\usepackage{fontspec}
\setmainfont{Latin Modern Roman}
\usepackage{amsmath,amssymb,mathtools}
\usepackage{unicode-math}
\setmathfont{Latin Modern Math}
\usepackage{enumitem}
\usepackage{xurl,hyperref,bookmark}
\hypersetup{colorlinks=true,urlcolor=blue}
\setcounter{secnumdepth}{0}
\setlength{\parindent}{0pt}
\setlength{\parskip}{5pt}
\setlength{\emergencystretch}{3em}
\begin{document}
\section*{Nearby Lagrangian conjecture}

\subsection{OpenAI's claim}
\textit{Nearby Lagrangian counterexample}: For some sufficiently large even $N$, the paper claims a closed embedded exact Lagrangian in $T^*(S^9\times S^{N-1})$ that cannot be taken to the zero section by a Hamiltonian isotopy, despite sharing the base’s diffeomorphism type. Thus the unrestricted nearby-Lagrangian assertion is contradicted.

\subsection{Mathematical statement of the claimed result}
\paragraph{Mathematical claim 1.} For some sufficiently large even $N$, let $Q=S^9\times S^{N-1}$. The manuscript claims a closed exact smoothly embedded Lagrangian $L\subset T^*Q$ with $L\cong Q$ as smooth manifolds, but with no compactly supported Hamiltonian isotopy sending the zero section onto $L$. Exactness means that the canonical cotangent one-form restricts to $dF$ on $L$ for some function $F$.

\subsection{Outline of the argument}
\textit{Nearby Lagrangian counterexample}: Generating families are organized by their homogeneous quadratic behavior at large vertical radius. Stable tube data determine spherical fibrations, and a nontrivial obstruction is selected from the tube-to-fibration comparison. A regular vertical-gradient construction produces an exact immersion, followed by a separate injectivity argument; the retained stable obstruction prevents a Hamiltonian isotopy to the zero section.

\subsection{Scope relative to this problem}
\textit{Nearby Lagrangian counterexample}: The constructed Lagrangian is closed, exact and smoothly embedded in a cotangent bundle, but is not Hamiltonian isotopic to its zero section; one permitted base disproves this universal assertion. Compact-support failure suffices because a Hamiltonian isotopy of compact Lagrangians has compact image and can be cut off without altering that image.

The claimed counterexample lies within the question’s hypotheses and contradicts its universal conclusion.

\subsection{Formal verification}
\textit{Nearby Lagrangian counterexample}: The pinned catalogue lists no Lean formalization for this result.

This site has not run the Lean development. These links describe OpenAI’s selected formal statements and should be read with their stated scope.

\subsection{Status caveat}
OpenAI reports that an internal model produced the manuscripts. Their arguments have not been independently verified as proofs by this site. Any portion outside the documented formal scope may contain mistakes, and the catalogue label reports the announced mathematical claim.

\subsection{Sources}
\begin{itemize}
\item OpenAI, \emph{A counterexample to the nearby Lagrangian conjecture}, 2026-09-23; citation key \texttt{OAI:A-counterexample-to-the-nearby-Lagrangian-conjecture-September-23-2026}. \href{https://github.com/openai/math/blob/adc7f1241b42e322a6451854ab7e4b4c146bf78a/preprints/A-counterexample-to-the-nearby-Lagrangian-conjecture-September-23-2026/paper.pdf}{PDF (GitHub)}; \href{https://github.com/openai/math/raw/adc7f1241b42e322a6451854ab7e4b4c146bf78a/preprints/A-counterexample-to-the-nearby-Lagrangian-conjecture-September-23-2026/paper.pdf}{PDF (direct)}; \href{https://github.com/openai/math/blob/adc7f1241b42e322a6451854ab7e4b4c146bf78a/preprints/A-counterexample-to-the-nearby-Lagrangian-conjecture-September-23-2026/README.md}{Manuscript page and citation}; \href{https://github.com/openai/math/blob/main/preprints/A-counterexample-to-the-nearby-Lagrangian-conjecture-September-23-2026/paper.pdf}{Latest version}; Statement source: \href{https://github.com/openai/math/blob/adc7f1241b42e322a6451854ab7e4b4c146bf78a/preprints/A-counterexample-to-the-nearby-Lagrangian-conjecture-September-23-2026/build/sections/introduction.tex\#L25}{thm:main}.
\item \href{https://github.com/openai/math/blob/adc7f1241b42e322a6451854ab7e4b4c146bf78a/CONTENTS.md}{OpenAI catalogue entry}.
\end{itemize}
\end{document}
